\documentclass[12pt]{beamer}
% \usetheme{umbc4}
\usetheme{moloch}

\usepackage{amsmath,amscd}
\usepackage{fontspec}
\usepackage{unicode-math}
\usepackage{pifont}
\usepackage{xcolor}
\usepackage[dvipsnames]{xcolor}
\usepackage{listings}
\usepackage{minted}
\usepackage{tikz}
\usepackage{twemojis}
\usepackage{nicematrix}

\usetikzlibrary{fit}

\setmonofont{DejaVu Sans Mono}

% aet the title bar to black background with white text
\setbeamercolor{frametitle}{bg=black, fg=white}

% ensure other structural elements (bullets, etc.) are visible
\setbeamercolor{structure}{fg=white}

\setmainfont{Libertinus Serif}
\setsansfont{Libertinus Sans}
\setmathfont{STIX Two Math}

% define colors for syntax highlighting
\definecolor{codegreen}{rgb}{0,0.6,0}
\definecolor{codegray}{rgb}{0.5,0.5,0.5}
\definecolor{codecoffee}{rgb}{0.75,0.6,0.3}
\definecolor{question}{rgb}{1.0,0.75,0.1}
\definecolor{lightblue}{rgb}{0.55,0.75,0.9}
\definecolor{yellow}{rgb}{1.0, 1.0, 0.0}
\definecolor{codecomment}{rgb}{0.97, 0.74, 0.4}
\definecolor{codepy}{rgb}{0.9, 0.85, 0.85} %{#F8BC65}
\definecolor{funcColor}{rgb}{0.47, 0.6, 0.8}
\definecolor{midgray}{rgb}{0.4, 0.4, 0.3}
\definecolor{indexwhite}{rgb}{0.9, 0.9, 0.9}

\setbeamercolor{background canvas}{bg=black}
\setbeamercolor{normal text}{fg=white}


\newcommand{\FrameTitle}[2]{\frametitle{#1 \hfill \small\textnormal{\textcolor{midgray}{#2}}}}
\newcommand{\snl}{\vspace*{8pt}}
\newcommand{\nl}{\vspace*{1em}\\}
\newcommand{\vsp}{\vspace*{1em}}

\newcommand{\Matrix}[1]{\textbf{#1}}

\begin{document}
\title{\center{Applied Linear Algebra \\\vsp Linear Functions \\}}

\author{\center{Devi Prasad M \\ Manipal School of Information Sciences \\ Manipal \\}}
\date{\center{Friday, 7 August 2026}}
\maketitle


\begin{frame}\FrameTitle{Linear Functions}{Unit 1}
    Consider the vector inner product function
    \[
        f(x) = a^{\textrm{T}} x = a_1x_1 + \cdots + a_nx_n
    \]
    where $a$ is a fixed $n$-vector, and $x$ is a variable $n$-vector.

    \snl The function $f$ is called a \emph{linear function} of $x$.

    \snl The vector $a$ is called the \emph{coefficient vector} of the linear function $f$.

    \snl $a$ is also called the \emph{weight vector} of the linear function $f$.
\end{frame}

\begin{frame}\FrameTitle{Linear Functions}{Unit 1}
    A linear function satisfies the following two properties:
    \begin{itemize}
        \item \emph{Additivity}. $f(x + y) = f(x) + f(y)$ for all $n$-vectors $x$ and $y$.
        \item \emph{Homogeneity}. $f(\alpha x) = \alpha\, f(x)$ for all $n$-vectors $x$ and all scalars $\alpha$.
    \end{itemize}

    \vsp The vector inner product function
    \[
        f(x) = a^{\textrm{T}} x = a_1x_1 + \cdots + a_nx_n, \ \ \  a: \mathbb{R}^n, \ \ x: \mathbb{R}^n
    \]
    is linear.
\end{frame}

\begin{frame}\FrameTitle{Superposition and Linearity}{Unit 1}
    \begin{definition}
        A function $f$ is \emph{linear} if it satisfies the \emph{superposition principle}:
        \[
            f(\alpha \, x + \beta \, y) = \alpha \, f(x) + \beta \,f(y)
        \]
        for all $n$-vectors $x$ and $y$, and all scalars $\alpha$ and $\beta$.
    \end{definition}

    \vsp
    Given the inner product function $f(x) = a^{\textrm{T}} x$, we can verify that it is linear:
    \[
        \begin{array}{rll}
        f(\alpha \, x + \beta \, y) & = & a^{\textrm{T}} (\alpha \, x + \beta \, y) \\
                                      & = & a^{\textrm{T}} (\alpha \, x) + a^{\textrm{T}} (\beta \, y) \\
                                      & = & \alpha (a^{\textrm{T}} \, x) + \beta (a^{\textrm{T}} \, y) \\
                                      & = & \alpha f(x) + \beta f(y)
        \end{array}
    \]
\end{frame}

\begin{frame}\FrameTitle{Linear Functions}{Unit 1}
\end{frame}

\begin{frame}\FrameTitle{Linear Functions}{Unit 1}
\end{frame}



\end{document}
